\section{Embedded Platform and Software Architecture}
\label{sec:embedded_platform}

The STM32 integration had three concurrent communication requirements: a
reliable \(1\,\mathrm{Mbit/s}\) servo bus, continuous E-Skin reception during
motion, and independent diagnostic output. Application code also had to remain
separated from STM32CubeMX-generated files. The implementation therefore uses
CubeMX as the hardware boundary and passes initialized peripheral handles to
the handwritten application modules.

\subsection{Microcontroller and Communication Configuration}

The firmware targets the NUCLEO-U545RE-Q. The MCU and LPUART1 clocks are derived
directly from the \(16\,\mathrm{MHz}\) high-speed internal oscillator. LPUART1
uses PA2 and PA3 with alternate function AF8 and operates bidirectionally at
\(1\,\mathrm{Mbit/s}\), 8-N-1, without hardware flow control. Selecting the
peripheral clock explicitly was necessary because the baud-rate setting alone
does not define the physical bit rate. During bring-up, increasing the system
clock from \(4\,\mathrm{MHz}\) to \(16\,\mathrm{MHz}\) and receiving each known
servo response as one block eliminated observed byte loss. The enabled RX FIFO
additionally buffers short bursts while another interrupt is serviced.

The large-area event-driven E-Skin follows the system described by Bergner,
Dean-Leon, and Cheng \cite{bergner2020eskin}. An ESP32-C6 sends its maximum
proximity value from GPIO16 to PC11/UART4\_RX at \(115200\,\mathrm{bit/s}\),
8-N-1. PC10/UART4\_TX is configured but unused by this unidirectional protocol.
A third UART connection, the ST-LINK virtual COM port at
\(115200\,\mathrm{bit/s}\), carries logs and controller telemetry. The three
peripherals shift data independently, although interrupts, parsing, and blocking
HAL calls still share processor time.

\subsection{Generated Code and UART Processing}

Clock, GPIO, interrupt, and peripheral initialization remain in the generated
\texttt{Core} directory, with the \texttt{.ioc} file as their configuration
source. The generated \texttt{main()} initializes LPUART1 and UART4, passes both
handles to \texttt{App\_Init()}, and repeatedly calls \texttt{App\_Process()}.
Handwritten UART, servo, kinematics, control, and sequencing code resides under
\texttt{App}. The only required post-generation adjustment, explicit servo RX
FIFO activation, remains inside a CubeMX user-code section.

Servo communication uses bounded blocking transfers because every request has a
known response length. E-Skin samples instead arrive continuously, so UART4
receives one byte at a time by interrupt. The callback ignores carriage returns
and publishes a frame only when a newline follows exactly five payload
characters, for example \texttt{0.123}. A sequence counter prevents foreground
code from reading a frame during an interrupt update. The conversion with
\texttt{strtof()} and the finite-value check run in
\texttt{UartCell\_Process()}, outside the interrupt. This keeps the interrupt
short and reduces interference with the faster servo bus. Counters for bytes,
valid and invalid frames, UART errors, receive-restart failures, and the last
valid sample support fault diagnosis; UART recovery is deferred to foreground
context.

\subsection{Application Integration and E-Skin Response}

\texttt{App\_Init()} starts E-Skin reception before any motion command and
requires a valid sample within \(2\,\mathrm{s}\). A missing sample or an initial
value of at least \(0.050\) causes a latched \texttt{FAULT}. Only after this check
does the application command the start pose and verify it with encoder feedback
and finite timeouts.

The demonstration calculates its initial tool center point through forward
kinematics and generates four waypoints in a constant-\(x\) plane. It exclusively
selects \texttt{Kinematics\_\allowbreak MoveEndEffector\allowbreak
ToPosition\allowbreak ResolvedRate()}, described in
Section~\ref{sec:feedback_control}. This blocking motion routine repeatedly
calls \texttt{app\_motion\_abort\_requested()}, which processes pending E-Skin
frames and the user button. Sensor supervision therefore remains active within
each control cycle rather than only between waypoints.

At an E-Skin value of \(0.050\), the application stops issuing trajectory
updates, reads the four active joint positions, and writes them back as hold
targets before entering \texttt{PAUSED\_SKIN}. Motion resumes at the same
waypoint only after the value remains at or below \(0.020\) for
\(500\,\mathrm{ms}\). A data age above \(1\,\mathrm{s}\), proximity during start
positioning, or a failed hold command instead produces \texttt{FAULT}. This is
only a software risk-reduction mechanism: response latency depends on the
current control and UART operation, the servos remain powered, and successful
holding still depends on communication. A physical means of disconnecting
actuator power therefore remains necessary.

\subsection{Control Integration and Individual Contribution}

The controller evaluation compares a one-shot baseline with Cartesian
resolved-rate feedback under the same geometry, limits, speed, acceleration,
and timeout. The baseline sends one complete inverse-kinematic solution and
uses encoder readings at detected joint standstill only to evaluate the
remaining model-based TCP error. Thus, ``without controller'' means without an
additional Cartesian loop; the internal STS3215 position controllers remain
active. The selected controller
synchronously reads all four active joints, reconstructs the TCP, calculates a
bounded Cartesian P step in the HAL-independent
\texttt{Control\_ResolvedRate\_CalculateStep()}, maps it through a damped
Jacobian inverse, and synchronously writes the updated targets. The optional
joint-tick PID remains in the source but is neither evaluated nor selected.

The main integration challenge was combining continuous low-rate sensor traffic
with tightly spaced servo responses and a blocking motion interface. Moving
floating-point parsing out of the UART4 interrupt, enabling the servo RX FIFO,
and reusing the motion abort callback solved this without an RTOS. Alessandro
Canalicchio's contribution comprised the STM32 and CubeMX configuration,
generated/application-code separation, diagnostics and fault propagation,
implementation and evaluation of the one-shot baseline through the common
kinematics and servo interfaces, integration of the resolved-rate controller
into the final application path, and the STM32-side E-Skin integration. Sensor
acquisition and preprocessing on the ESP32 belong to the separate sensing
subsystem.
